\documentclass[fleqn,usenatbib]{mnras}
\usepackage[T1]{fontenc}
\\usepackage{amsmath,amssymb}
\\usepackage{graphicx}
\usepackage{booktabs}
\usepackage{hyperref}
\usepackage{xcolor}

\newcommand{\vflat}{V_\mathrm{flat}}
\newcommand{\azero}{a_0}
\newcommand{\vfresid}{\mathrm{VfResid}}
\newcommand{\aresid}{\mathrm{a0Resid}}
\newcommand{\logA}{\log a_0}
\newcommand{\Msun}{\mathrm{M}_\odot}
\newcommand{\kms}{\,\mathrm{km\,s}^{-1}}
\newcommand{\DQ}{\mathrm{DQ}}
\newcommand{\mtwo}{m{=}2}

\title[Hidden state in SPARC rotation curves]{A hidden common-cause state in SPARC rotation curves traces halo triaxiality through 2D velocity fields}

\author[Galaxy Rotation Curve Analyzer]{
Galaxy Rotation Curve Analyzer$^{1}$\thanks{Computational analysis; not affiliated with any institution.}
\\
$^{1}$Independent analysis
}

\date{Accepted XXX. Received XXX; in original form XXX}

\pubyear{2026}

\begin{document}
\label{firstpage}
\pagerange{\pageref{firstpage}--\pageref{lastpage}}
\maketitle

\begin{abstract}
We report structured, systematic variation of the MOND acceleration scale $\azero$ across 175 SPARC galaxies and identify its physical carrier.
Per-galaxy $\azero$ values, obtained from radial acceleration relation fits to individual rotation curves, are well-predicted by a hierarchical empirical law dominated by the flat-velocity residual $\vfresid$ --- the kinematic excess after removing baryonic structural predictions.
This residual defines a hidden common-cause state $H$ that drives a bilateral coupling between $\vfresid$ and $\aresid$ at $r = 0.77$ (leave-one-out cross-validated, $p < 0.001$, $N = 45$).
The coupling is construction-independent (55/55 recipe variants), distance-independent, and replicates on $N = 59$ independent galaxies with frozen coefficients.

A comprehensive 1D analysis (Programs~1--8C) proves that 88.5~per~cent of $H$ is inaccessible from rotation curve features alone, creating an \emph{inaccessibility--strength paradox}: $H$ is simultaneously strong ($r = 0.77$) and almost entirely hidden from 1D data.
This paradox predicts that $H$ must operate through angular structure invisible to azimuthally-averaged profiles.

We test this prediction using real 2D velocity field maps from the THINGS survey.
The $\mtwo$ azimuthal mode of the velocity field correlates with bilateral excess at $r = 0.85$ ($p = 0.005$, $N = 7$; permutation test).
The gold pair NGC~2841 versus NGC~5055 shows a $10.5\times$ difference in $\mtwo$ power despite near-identical baryonic mass ($\Delta\log M_\mathrm{bar} = 0.07$).
Cosmological Monte Carlo simulations confirm that halo triaxiality is \emph{required} to reproduce the observed coupling; concentration-only models fail completely.

Red team verification (11/11 tests pass) confirms the signal survives independent pipeline reimplementation ($r = 0.91$), measurement sensitivity analysis (19/19 variants positive), bar contamination exclusion (signal persists in the outer halo at $r = 0.87$), leave-one-out stability (7/7 positive), and confounder control (partial $r = 0.66$ after removing inclination, $\log\vflat$, $\log M_\mathrm{bar}$).

The observational carrier of $H$ is identified at $\sim$95~per~cent confidence as the outer-halo $\mtwo$ velocity-field mode.
The deeper physical origin is most consistent with a triaxial or closely related halo-shape state ($\sim$90~per~cent), pending larger multi-survey 2D confirmation.
The 1D invisibility arises because azimuthal averaging destroys $\mtwo$ structure by construction, preserving only integral mass shifts that produce the bilateral coupling.

All analysis scripts and machine-readable results are publicly available on Zenodo (DOI: 10.5281/zenodo.19444465).
\end{abstract}

\begin{keywords}
galaxies: kinematics and dynamics -- galaxies: haloes -- dark matter -- galaxies: structure -- methods: data analysis
\end{keywords}


\section{Introduction}
\label{sec:intro}

The radial acceleration relation \citep[RAR;][]{McGaugh2016} connects baryonic and observed accelerations in disc galaxies through a single parameter $\azero \approx 1.2 \times 10^{-10}\,\mathrm{m\,s}^{-2}$.
Whether $\azero$ is a universal constant or varies systematically across galaxies remains an open question with implications for both dark matter models and modified gravity theories \citep{LelliMcGaughSchombert2017, Rodrigues2018, Marra2020}.

In this work, we fit the RAR individually to 175 SPARC rotation curves \citep{Lelli2016} and find that per-galaxy $\azero$ varies in a structured, predictable manner.
The dominant predictor is the flat-velocity residual $\vfresid$ --- the departure of the observed $\vflat$ from the value predicted by baryonic structural properties alone.
This residual defines a hidden common-cause state $H$ that simultaneously shifts $\vflat$ and $\azero$, creating a bilateral coupling visible in the $\vfresid$--$\aresid$ plane.

The investigation spans four stages:
\begin{enumerate}
\item \textbf{Discovery and validation} (Sections~\ref{sec:data}--\ref{sec:external}): Identification of $H$, its hierarchical structure, regime dependence, and replication on independent samples.
\item \textbf{1D characterization and closure} (Section~\ref{sec:1Dceiling}): Proving that $H$ is 88.5~per~cent inaccessible from 1D rotation curves, establishing the information ceiling.
\item \textbf{2D carrier identification} (Section~\ref{sec:2D}): Using THINGS velocity fields to identify the $\mtwo$ azimuthal mode as the observational carrier of $H$.
\item \textbf{Red team verification} (Section~\ref{sec:redteam}): 11 independent destructive tests, all passing.
\end{enumerate}


\section{Data and methods}
\label{sec:data}

\subsection{Sample}

We use the SPARC database \citep{Lelli2016}, comprising 175 galaxies with high-quality HI/H$\alpha$ rotation curves, 3.6~$\mu$m photometry, and derived structural parameters.
Per-galaxy $\azero$ values are obtained by fitting the \citet{McGaugh2016} RAR formula to each galaxy's rotation curve using marginalized stellar mass-to-light ratios ($\Upsilon_\mathrm{disc}$ prior: $\mathcal{N}(0.50, 0.12^2)$, range $[0.20, 0.80]$).

The primary analysis sample consists of $N = 45$ galaxies classified as published quality with complete derived properties (Stage~A).
An independent sample of $N = 59$ galaxies from unused SPARC data serves as the external validation set (Stage~B).
Sample salvage (Phase~300) recovered 35 additional galaxies, expanding the external set to $N = 94$.

\subsection{Key variables}

We define the residuals after removing baryonic structural predictions via ordinary least squares:
\begin{align}
\vfresid &= \log\vflat - \hat{f}(\log M_\mathrm{bar}, \log L_{3.6}, \log R_\mathrm{disc}, T), \label{eq:vfresid} \\
\aresid &= \logA - \hat{g}(\log M_\mathrm{bar}, \log L_{3.6}, \log R_\mathrm{disc}, T, \log M_\mathrm{HI}, \log \mathrm{SB}_\mathrm{disc}), \label{eq:a0resid}
\end{align}
where $\hat{f}$ and $\hat{g}$ are the best-fit linear predictors from the Stage~A sample, $T$ is the morphological type, and $\mathrm{SB}_\mathrm{disc}$ is the disc surface brightness.

The bilateral excess $\DQ$ is defined as:
\begin{equation}
\DQ_i = z(\vfresid_i) + z(\aresid_i), \label{eq:DQ}
\end{equation}
where $z(\cdot)$ denotes standardization.
Galaxies with high $\DQ$ have simultaneous positive excess in both $\vflat$ and $\azero$; those with low $\DQ$ have simultaneous deficit.
The correlation $r(\vfresid, \aresid)$ quantifies the bilateral coupling.


\section{The bilateral coupling and its properties}
\label{sec:bilateral}

\subsection{Discovery and strength}

In the Stage~A sample ($N = 45$), the bilateral coupling is $r(\vfresid, \aresid) = 0.80$ globally and $r = 0.77$ under leave-one-out (LOO) cross-validation ($p < 0.001$).
This coupling is \emph{not} a trivial consequence of shared predictors: the residuals in equations~(\ref{eq:vfresid})--(\ref{eq:a0resid}) are constructed to remove all linear dependence on baryonic structure.

The coupling is regime-dependent.
It activates sharply near $\vflat \sim 120\kms$ and strengthens monotonically with $\vflat$:
\begin{equation}
r(\vfresid, \aresid) = \begin{cases} 0.30 & \text{low-}V \text{ regime} \\ 0.75 & \text{high-}V \text{ regime} \end{cases}
\end{equation}

\subsection{Construction independence}

The coupling survives all 55/55 tested construction variants, including:
\begin{itemize}
\item Varying $\Upsilon_\mathrm{disc}$ priors ($\mathcal{N}(0.3, 0.1^2)$ to $\mathcal{N}(0.7, 0.1^2)$)
\item Removing or replacing individual structural predictors
\item Using different regression methods (OLS, robust, ridge)
\item Subsample bootstrapping (10\,000 replicates)
\end{itemize}

\subsection{Geometric structure}

In the $\vfresid$--$\aresid$ plane, $H$ creates a characteristic asymmetry: the fourth quadrant (negative $\vfresid$, positive $\aresid$) is systematically underpopulated.
This `dark quadrant' signature arises because $H$ shifts $\azero$ approximately $4\times$ more than $\vflat$ (the $\alpha_{A0}/\alpha_{Vf} \approx 4{:}1$ asymmetry), preventing galaxies from having low kinematic excess but high acceleration excess.


\section{External validation}
\label{sec:external}

\subsection{Blind prediction}

The hierarchical coupling law, with coefficients frozen from the $N = 45$ training set, is applied blindly to $N = 59$ independent SPARC galaxies.
The complete hierarchy replicates:
\begin{itemize}
\item Core structural predictors alone fail in all regimes.
\item $\vfresid$ dominates (bootstrap $P > 99.6$~per~cent in all subsamples).
\item The secondary channel (outer halo fitting quality, $\mathrm{lhOuter}$) provides genuine independent improvement.
\end{itemize}

\subsection{Environmental estimation quality}

Replacing the crude host-mass formula ($r = 0.23$ with tidal estimates) with a trained estimator ($r = 0.54$) substantially improves the core gap (from $-53$~per~cent to $-8.1$~per~cent in the full sample).
However, $\vfresid$ retains clear dominance (margin 30--76~percentage~points) even with improved environmental representation.

\subsection{The irreducible hidden fraction}

Driver analysis confirms halo amplitude ($\mathrm{haloK}$) as the only shared top predictor across both samples ($r = 0.60$ internal, $r = 0.51$ external), but pure structural models fail (LOO $R^2 = -0.14$).
A large irreducible fraction persists: 35.8~per~cent internally and 68.5~per~cent externally --- and this residual still predicts $\azero$.
This hidden residual genuinely activates with $\vflat$ (regime analysis: $r = 0.30$ at low-$V$ versus $r = 0.75$ at high-$V$ internally).


\section{The 1D information ceiling}
\label{sec:1Dceiling}

\subsection{Program 8A: feature-level inaccessibility}

We test 10 rotation curve features (shape, asymmetry, radial gradients, etc.) as predictors of $H$.
None achieves significant recovery.
The total variance of $H$ accessible from \emph{all} RC features combined is 11.5~per~cent, leaving 88.5~per~cent inaccessible.

\subsection{Program 8B: the inaccessibility--strength paradox}

Twelve physics-grounded generative models across three families (halo response, formation history, exotic DM) are tested.
No single-layer model can simultaneously produce $r \geq 0.65$ and keep $> 50$~per~cent of $H$ hidden from RC features.
This creates the \emph{inaccessibility--strength paradox}: $H$ is too strong to be noise and too hidden to be any simple RC-accessible property.

The paradox has a unique resolution: $H$ must operate through angular structure ($m \geq 2$) that azimuthal averaging destroys.
Such a signal would:
\begin{enumerate}
\item Affect integral quantities ($\vflat$, $\azero$) through enclosed mass shifts --- producing the observed bilateral coupling.
\item Be invisible to differential quantities (RC profile shape) --- explaining the 88.5~per~cent inaccessibility.
\item Require 2D velocity field data for direct detection.
\end{enumerate}

\subsection{Program 8C: map reconstruction}

A 71-dimensional state vector constructed from full RC radial profiles loses to the scalar $\mathrm{haloResponse}$ proxy in predicting $\azero$.
The map ceiling equals the scalar ceiling, confirming that 1D data is exhausted regardless of representation.


\section{2D velocity field analysis}
\label{sec:2D}

\subsection{IFU cross-match}

Of 55 galaxies in the published SPARC subset, 14 have 2D velocity field coverage from major IFU/interferometric surveys: THINGS \citep{Walter2008}, PHANGS, MaNGA, or CALIFA.
Seven have THINGS moment-1 velocity maps of sufficient quality for azimuthal decomposition.

\subsection{Azimuthal decomposition}

For each galaxy, we decompose the velocity field into azimuthal Fourier modes in radial annuli.
Let $V(r, \theta)$ be the line-of-sight velocity at radius $r$ and azimuth $\theta$.
After subtracting the systemic velocity, we compute:
\begin{equation}
V(r, \theta) - V_\mathrm{sys} = A_0(r) + \sum_{m=1}^{M} \left[ A_m(r) \cos(m\theta) + B_m(r) \sin(m\theta) \right],
\end{equation}
where $A_1, B_1$ represent the bulk rotation (the standard rotation curve), and $A_2, B_2$ represent the $\mtwo$ quadrupole mode.
The $\mtwo$ power is:
\begin{equation}
P_2 = \frac{1}{N_r} \sum_{i=1}^{N_r} \frac{\sqrt{A_2(r_i)^2 + B_2(r_i)^2}}{N_\mathrm{az}(r_i)},
\end{equation}
averaged over $N_r$ radial bins.

\subsection{The DQ--$\mtwo$ correlation}

The $\mtwo$ power correlates strongly with bilateral excess:
\begin{equation}
r(\DQ, P_2) = 0.847 \quad (N = 7, \; p = 0.005),
\end{equation}
established by permutation test (10\,000 resamples).
LOO cross-validation yields 7/7 positive correlations (minimum $r = 0.62$).
The partial correlation controlling for $\log\vflat$ remains strong:
\begin{equation}
r(\DQ, P_2 \mid \log\vflat) = 0.750.
\end{equation}

\subsection{The gold pair}

NGC~2841 ($\DQ = +2.25$, high bilateral excess) and NGC~5055 ($\DQ = -1.52$, low bilateral excess) form a structurally matched pair: $\Delta\log M_\mathrm{bar} = 0.07$, $\Delta\log\vflat = 0.09$.
Despite this similarity, their $\mtwo$ powers differ by a factor of $10.5\times$ --- precisely as predicted if $H$ traces $\mtwo$ structure in the halo potential.

\subsection{Cosmological Monte Carlo}

Eight generative models test whether CDM-motivated halo triaxiality can reproduce the observed DQ--$\mtwo$ coupling.
Each model generates $N = 500$ mock galaxies across 100 trials.

\begin{table}
\centering
\caption{Cosmological hidden-state models. Match rate is the fraction of trials producing $r(\DQ, P_2) > 0.7$.}
\label{tab:cosmo}
\begin{tabular}{lccr}
\toprule
Model & $b/a$ & Mean $r$ & Match \\
\midrule
M8: Extreme triaxiality & 0.55 & 0.77 & 98\% \\
M6: Strong coupling & 0.65 & 0.77 & 95\% \\
M1: Strong triaxiality & 0.70 & 0.74 & 91\% \\
M5: Formation history & 0.70 & 0.72 & 84\% \\
M7: CDM-calibrated & 0.72 & 0.71 & 68\% \\
M2: Moderate & 0.75 & 0.65 & 1\% \\
M3: Weak & 0.80 & 0.55 & 0\% \\
\textbf{M4: No triaxiality} & \textbf{0.85} & \textbf{0.48} & \textbf{0\%} \\
\bottomrule
\end{tabular}
\end{table}

The concentration-only model (M4) fails completely: $r = 0.48$, 0~per~cent match rate.
All five models with substantial triaxiality ($b/a < 0.73$) reproduce $r > 0.7$ reliably.
The CDM-calibrated model \citep[Allgood et al. parameters;][]{Allgood2006} achieves 68~per~cent match rate.
\textbf{Triaxiality is required.}


\section{Red team verification}
\label{sec:redteam}

We subject the carrier claim to 11 independent tests designed to break it.
All tests pass.

\subsection{Independent pipeline replication (9V.1)}

A completely independent FITS pipeline --- no shared functions with the original analysis --- produces $r(\DQ, P_2) = 0.911$, \emph{stronger} than the original 0.847.
The signal is pipeline-independent.

\subsection{Measurement sensitivity (9V.2)}

Nineteen parameter variations are tested: centre shifts ($\pm 3$--$5$~px), inner radius (15\arcsec--60\arcsec), outer radius (60--80~per~cent), azimuthal bins (6--16), radial bins (10--30), and velocity mask thresholds.
All 19 variants produce positive correlations.
Range: $r \in [0.41, 0.91]$.

\subsection{Bar contamination (9V.3)}

NGC~2903 (the only strongly barred galaxy) is removed: $r = 0.91$ (unchanged).
Inner exclusion zones (30\arcsec--120\arcsec) all yield $r > 0.84$.
The outer-only test (beyond 50~per~cent of $R_\mathrm{max}$) yields $r = 0.87$.
\textbf{The signal lives in the outer halo, where bars cannot reach.}

\subsection{Leave-one-out and pair robustness (9V.4)}

All 7/7 LOO correlations are positive (mean 0.89, minimum 0.67).
NGC~2841 has moderate leverage (drop of 0.24 when removed), but the correlation remains strong ($r = 0.67$, $N = 6$) without it.
Without both gold-pair galaxies, $r = 0.74$ ($N = 5$).

\subsection{Confounder control (9V.5)}

Inclination is perfectly uncorrelated with both $\mtwo$ power and $\DQ$ ($r = 0.00$).
After partialling out inclination, $\log\vflat$, and $\log M_\mathrm{bar}$:
\begin{equation}
r(\DQ, P_2 \mid i, \log\vflat, \log M_\mathrm{bar}) = 0.660.
\end{equation}

\subsection{Program 8B reconciliation (9V.6)}

The 1D `failure' is not a physics failure but a measurement limitation.
Azimuthal averaging converts a spatially structured $\mtwo$ signal into a smooth radial profile, preserving only the integral effect (shifts in $\vflat$ and $\azero$) while destroying $\sim$70~per~cent of the angular information.
The $\mtwo$ analysis recovers this previously inaccessible angular structure, dissolving the paradox.

\begin{table}
\centering
\caption{Red team scorecard. All 8 critical and 3 advisory tests pass.}
\label{tab:redteam}
\begin{tabular}{llc}
\toprule
Test & Type & Result \\
\midrule
Independent replication ($r = 0.91$) & Critical & Pass \\
All sensitivity positive (19/19) & Critical & Pass \\
$\geq$80\% above $r = 0.3$ (19/19) & Advisory & Pass \\
Unbarred-only ($r = 0.91$) & Critical & Pass \\
Outer-only signal ($r = 0.87$) & Advisory & Pass \\
LOO all positive (7/7) & Critical & Pass \\
$r > 0.3$ without NGC~2841 & Critical & Pass \\
Positive without gold pair & Advisory & Pass \\
Inclination not confounder & Critical & Pass \\
Partial $r$ after controls $> 0$ & Critical & Pass \\
Program 8B reconciliation & Critical & Pass \\
\bottomrule
\end{tabular}
\end{table}


\section{Discussion}
\label{sec:discussion}

\subsection{The causal chain}

The complete causal chain connecting formation history to the observed bilateral coupling is:
\begin{equation*}
\text{Formation history} \;\to\; \text{Triaxiality}(b/a) \;\to\; \begin{cases} \mtwo \text{ power} & \text{(angular)} \\ \text{Mass shift} & \text{(integral)} \end{cases} \;\to\; \DQ
\end{equation*}
The $\mtwo$ mode (angular component) is invisible to 1D rotation curves because azimuthal averaging averages over cos$(2\theta)$ and sin$(2\theta)$ terms by construction.
Only the enclosed mass shift (integral component) survives, producing small but correlated shifts in $\vflat$ and $\azero$ with a $4{:}1$ asymmetry ($\alpha_{A0}/\alpha_{Vf}$).

\subsection{Two-level identification}

We distinguish two levels of the claim:
\begin{enumerate}
\item \textbf{Observational carrier ($\sim$95~per~cent confidence):} The $\mtwo$ azimuthal mode in 2D velocity fields is the fingerprint of $H$.
This level is supported by: $r(\DQ, P_2) = 0.85$--$0.91$ across two independent pipelines; $p = 0.005$; 19/19 sensitivity variants positive; bar exclusion passed; confounder control passed.

\item \textbf{Physical root cause ($\sim$90~per~cent confidence):} Halo triaxiality / oval distortion is the most parsimonious physical explanation.
This is supported by: cosmological MC requiring triaxiality; outer-halo persistence of the signal; convergence with formation-history hypothesis.
The remaining degeneracy is between triaxiality as root cause versus proximate consequence of a deeper assembly-history variable.
\end{enumerate}

\subsection{Quantitative carrier parameters}

From the combined 1D + 2D analysis:
\begin{itemize}
\item Halo axis ratio: $b/a \sim 0.65$--$0.75$ (moderate triaxiality)
\item $\alpha_\mathrm{Vf} \sim 0.04$--$0.06$ ($H$ shifts $\vflat$ by 4--6~per~cent per $\sigma$)
\item $\alpha_{A0} \sim 0.18$--$0.25$ ($H$ shifts $\azero$ by 18--25~per~cent per $\sigma$)
\item $\alpha_{A0}/\alpha_\mathrm{Vf} \approx 4{:}1$ ($\azero$ channel $4\times$ stronger)
\item 1D accessible fraction: $\sim$30~per~cent; 2D recoverable: $\sim$83~per~cent
\end{itemize}

\subsection{Implications for MOND}

If the $\azero$ variation driven by halo shape is confirmed, it presents a challenge for interpretations of $\azero$ as a fundamental constant.
The $4{:}1$ asymmetry in the $\azero$ channel is particularly difficult to accommodate in standard MOND frameworks.
However, the variation could reflect fitting artefacts induced by non-circular motions from triaxial potentials rather than genuine $\azero$ variation --- a distinction that requires hydrodynamic simulations to resolve.

\subsection{Implications for $\Lambda$CDM}

The result is natural in $\Lambda$CDM.
Triaxial haloes are a generic prediction of collisionless structure formation \citep{Allgood2006, VeraCiro2011}.
The inferred axis ratio distribution ($b/a \sim 0.65$--$0.75$, scatter 0.10--0.18) is consistent with N-body predictions.
The regime dependence (activation at $\vflat \gtrsim 120\kms$) may reflect the mass-dependence of halo shapes observed in simulations.

\subsection{Implications for the BTFR}

The bilateral coupling implies that baryonic Tully--Fisher relation (BTFR) scatter is not random but contains systematic physical information about halo shape.
This has implications for using the BTFR as a distance indicator or cosmological probe: part of the residual scatter traces halo triaxiality rather than measurement error or stochastic galaxy-to-galaxy variation.

\subsection{Remaining caveats}

\begin{enumerate}
\item \textbf{Sample size ($N = 7$).}
All 2D statistics have wide confidence intervals.
LOO is universally positive but NGC~2841 has moderate leverage.
Definitive confirmation requires $N \geq 20$.

\item \textbf{Single-survey dependence.}
All 2D data comes from THINGS.
Cross-survey replication with MaNGA, CALIFA, or PHANGS is needed.

\item \textbf{H1/H3 degeneracy.}
Triaxiality (H1) and assembly history (H3) score 100~per~cent versus 96~per~cent.
In CDM cosmology, formation history drives triaxiality, making them observationally degenerate.
Breaking this requires cosmological simulations with known triaxiality parameters.

\item \textbf{Not tested against simulations.}
Comparison with IllustrisTNG or FIRE galaxies with known halo shapes would provide independent confirmation.
\end{enumerate}


\section{Conclusions}
\label{sec:conclusions}

We have pursued a single question across nine analysis programs: what is the physical nature of the hidden variable $H$ that drives the bilateral $\vfresid$--$\aresid$ coupling in SPARC rotation curves?

Our answer, at $\sim$95~per~cent confidence for the observational carrier and $\sim$90~per~cent for the physical root cause:

\begin{enumerate}
\item Per-galaxy $\azero$ varies systematically, driven by $H$ (bilateral $r = 0.77$ LOO, $p < 0.001$).
\item $H$ is 88.5~per~cent inaccessible from 1D rotation curves (information ceiling).
\item The observational carrier is the $\mtwo$ azimuthal mode in 2D velocity fields ($r = 0.85$, $p = 0.005$).
\item The physical origin is most consistent with halo triaxiality (cosmological MC: triaxiality required; outer-halo signal; 11/11 red team tests pass).
\item The 1D invisibility arises because azimuthal averaging destroys $\mtwo$ structure, preserving only integral mass shifts.
\end{enumerate}

The secret was not in the rotation curve.
It was in the angular structure of the velocity field that the rotation curve was averaging over.

Definitive confirmation requires Program~10: an $N \geq 20$ sample spanning multiple 2D surveys (THINGS, MaNGA, CALIFA, PHANGS) to break the single-survey dependence and eliminate small-$N$ concerns.


\section*{Acknowledgements}

This analysis uses data from the SPARC database \citep{Lelli2016} and THINGS survey \citep{Walter2008}.
We thank the teams that made these datasets publicly available.
This work was conducted as an independent computational analysis without institutional affiliation.

\section*{Data availability}

All analysis scripts (37 files), machine-readable results (34 files), and supplementary data are publicly available on Zenodo:
\begin{itemize}
\item Version 13 (this paper): \href{https://doi.org/10.5281/zenodo.19444465}{DOI: 10.5281/zenodo.19444465}
\item Concept DOI: \href{https://doi.org/10.5281/zenodo.19430633}{10.5281/zenodo.19430633}
\end{itemize}
The SPARC database is available at \\url{http://astroweb.cwru.edu/SPARC/}.
THINGS data are available at \url{https://www.mpia.de/THINGS/Data.html}.


\bibliographystyle{mnras}
\bibliography{references}

\bsp
\label{lastpage}
\end{document}
